\section{Conclusion} \label{sec:conclusion}

Ce projet a permis de définir et tester une mesure de l'efficacité adaptée à un graphe dynamique. 
Il a également permis de définir des modélisations de manifestations comme des attaques 
sur un graphe dynamique.
Nous avons pû apporter différentes approches pour calculer ces mesures, car la plupart de la littérature s'intéresse aux graphes dynamiques qui
changent en temps réel.

Les résultats obtenus sont prometteurs et mettent en évidence l'influence de l'aspect 
dynamique du graphe sur l'efficacité des stratégies d'attaque.
Cependant, nous avons vu que cela dépend fortement des choix de modélisation des attaques.

Nous pourrions changer la modélisation de notre problème, ajuster 
la mesure de l'efficacité ou encore tester d'autres stratégies d'attaque 
pour obtenir de meilleurs résultats. 

Cependant, ce projet a permis de mettre en lumière le fait qu'il était 
nécessaire de créer et tester des mesures adaptées aux graphes dynamiques 
afin de pouvoir utiliser ceux-ci pour modéliser des problèmes réels.